\documentclass[10pt,a4paper]{amsart}
\usepackage{amsmath,amssymb,amsthm,graphicx}
% Retained minimal preamble for the Stacks Project chapter varieties.tex.
% Verbatim excerpt of the chapter preamble: Xy-pic setup (native preamble lines 14-18), the native theorem
% declarations with the mathematical macros (native preamble lines 154-217), and the project own no-op
% definitions of the reference/slogan/history display environments (native wrapper.tex lines 1-3) that the
% counted statement uses for its slogan. Copied mechanically; nothing added, removed or redefined.
% For commutative diagrams we use Xy-pic
\usepackage[all]{xy}

% We use 2cell for 2-commutative diagrams.
\xyoption{2cell}
% Theorem environments.
%
\theoremstyle{plain}
\newtheorem{theorem}[subsection]{Theorem}
\newtheorem{proposition}[subsection]{Proposition}
\newtheorem{lemma}[subsection]{Lemma}

\theoremstyle{definition}
\newtheorem{definition}[subsection]{Definition}
\newtheorem{example}[subsection]{Example}
\newtheorem{exercise}[subsection]{Exercise}
\newtheorem{situation}[subsection]{Situation}

\theoremstyle{remark}
\newtheorem{remark}[subsection]{Remark}
\newtheorem{remarks}[subsection]{Remarks}

\numberwithin{equation}{subsection}

% Macros
%
\def\lim{\mathop{\mathrm{lim}}\nolimits}
\def\colim{\mathop{\mathrm{colim}}\nolimits}
\def\Spec{\mathop{\mathrm{Spec}}}
\def\Hom{\mathop{\mathrm{Hom}}\nolimits}
\def\Ext{\mathop{\mathrm{Ext}}\nolimits}
\def\SheafHom{\mathop{\mathcal{H}\!\mathit{om}}\nolimits}
\def\SheafExt{\mathop{\mathcal{E}\!\mathit{xt}}\nolimits}
\def\Sch{\mathit{Sch}}
\def\Mor{\mathop{\mathrm{Mor}}\nolimits}
\def\Ob{\mathop{\mathrm{Ob}}\nolimits}
\def\Sh{\mathop{\mathit{Sh}}\nolimits}
\def\NL{\mathop{N\!L}\nolimits}
\def\CH{\mathop{\mathrm{CH}}\nolimits}
\def\proetale{{pro\text{-}\acute{e}tale}}
\def\etale{{\acute{e}tale}}
\def\QCoh{\mathit{QCoh}}
\def\Ker{\mathop{\mathrm{Ker}}}
\def\Im{\mathop{\mathrm{Im}}}
\def\Coker{\mathop{\mathrm{Coker}}}
\def\Coim{\mathop{\mathrm{Coim}}}

% Boxtimes
%
\DeclareMathSymbol{\boxtimes}{\mathbin}{AMSa}{"02}

%
% Macros for moduli stacks/spaces
%
\def\QCohstack{\mathcal{QC}\!\mathit{oh}}
\def\Cohstack{\mathcal{C}\!\mathit{oh}}
\def\Spacesstack{\mathcal{S}\!\mathit{paces}}
\def\Quotfunctor{\mathrm{Quot}}
\def\Hilbfunctor{\mathrm{Hilb}}
\def\Curvesstack{\mathcal{C}\!\mathit{urves}}
\def\Polarizedstack{\mathcal{P}\!\mathit{olarized}}
\def\Complexesstack{\mathcal{C}\!\mathit{omplexes}}
% \Pic is the operator that assigns to X its picard group, usage \Pic(X)
% \Picardstack_{X/B} denotes the Picard stack of X over B
% \Picardfunctor_{X/B} denotes the Picard functor of X over B
\def\Pic{\mathop{\mathrm{Pic}}\nolimits}
\def\Picardstack{\mathcal{P}\!\mathit{ic}}
\def\Picardfunctor{\mathrm{Pic}}
\def\Deformationcategory{\mathcal{D}\!\mathit{ef}}
\newenvironment{reference}{}{}
\newenvironment{slogan}{}{}
\newenvironment{history}{}{}
\usepackage{cleveref}
\crefname{originalreference}{original reference}{original references}
\Crefname{originalreference}{Original reference}{Original references}
\makeatletter
\newlabel{morphisms-definition-degree}{{02NY}{source}{}{sourcepointer-1}{}}
\newlabel{morphisms-definition-degree@cref}{{[originalreference][0][]02NY}{[1][0][]source}{}{}{}}
\makeatother
\crefname{originallemma}{original lemma}{original lemmas}
\Crefname{originallemma}{Original lemma}{Original lemmas}
\makeatletter
\newlabel{lemma-euler-characteristic-morphism}{{0BEK}{source}{}{sourcepointer-2}{}}
\newlabel{lemma-euler-characteristic-morphism@cref}{{[originallemma][0][]0BEK}{[1][0][]source}{}{}{}}
\makeatother
\crefname{originalreference}{original reference}{original references}
\Crefname{originalreference}{Original reference}{Original references}
\makeatletter
\newlabel{cohomology-lemma-projection-formula}{{01E8}{source}{}{sourcepointer-3}{}}
\newlabel{cohomology-lemma-projection-formula@cref}{{[originalreference][0][]01E8}{[1][0][]source}{}{}{}}
\makeatother
\crefname{originallemma}{original lemma}{original lemmas}
\Crefname{originallemma}{Original lemma}{Original lemmas}
\makeatletter
\newlabel{lemma-numerical-polynomial-from-euler}{{0BEM}{source}{}{sourcepointer-4}{}}
\newlabel{lemma-numerical-polynomial-from-euler@cref}{{[originallemma][0][]0BEM}{[1][0][]source}{}{}{}}
\makeatother
\crefname{originallemma}{original lemma}{original lemmas}
\Crefname{originallemma}{Original lemma}{Original lemmas}
\makeatletter
\newlabel{lemma-dimension-locally-algebraic}{{0A21}{source}{}{sourcepointer-5}{}}
\newlabel{lemma-dimension-locally-algebraic@cref}{{[originallemma][0][]0A21}{[1][0][]source}{}{}{}}
\makeatother
\crefname{originallemma}{original lemma}{original lemmas}
\Crefname{originallemma}{Original lemma}{Original lemmas}
\makeatletter
\newlabel{lemma-dimension-fibres-locally-algebraic}{{0B2L}{source}{}{sourcepointer-6}{}}
\newlabel{lemma-dimension-fibres-locally-algebraic@cref}{{[originallemma][0][]0B2L}{[1][0][]source}{}{}{}}
\makeatother
\crefname{originalreference}{original reference}{original references}
\Crefname{originalreference}{Original reference}{Original references}
\makeatletter
\newlabel{morphisms-lemma-generically-finite}{{02NW}{source}{}{sourcepointer-7}{}}
\newlabel{morphisms-lemma-generically-finite@cref}{{[originalreference][0][]02NW}{[1][0][]source}{}{}{}}
\makeatother
\crefname{originalreference}{original reference}{original references}
\Crefname{originalreference}{Original reference}{Original references}
\makeatletter
\newlabel{coherent-lemma-finite-pushforward-coherent}{{01Y6}{source}{}{sourcepointer-8}{}}
\newlabel{coherent-lemma-finite-pushforward-coherent@cref}{{[originalreference][0][]01Y6}{[1][0][]source}{}{}{}}
\makeatother
\crefname{originallemma}{original lemma}{original lemmas}
\Crefname{originallemma}{Original lemma}{Original lemmas}
\makeatletter
\newlabel{lemma-numerical-polynomial-leading-term}{{0BEN}{source}{}{sourcepointer-9}{}}
\newlabel{lemma-numerical-polynomial-leading-term@cref}{{[originallemma][0][]0BEN}{[1][0][]source}{}{}{}}
\makeatother
\title{Intersection numbers of pulled-back invertible sheaves scale by the degree of the morphism restricted to the subscheme}
\author{The Stacks Project Authors}
\date{stacks-project:varieties.tex:local-snapshot-20261011; GNU FDL 1.2 or later}
\begin{document}
\maketitle
\noindent\textit{Editorial note.} All delivered mathematical text is the author's own, copied verbatim from the pinned Stacks Project chapter file varieties.tex; no statement or proof step was rewritten, reordered or removed.\par\medskip
\noindent\textit{Editorial note.} The counted span cites nine cross-chapter labels (including a definition), each recorded in assembly.external\_references with its canonical tag as a navigation pointer.\par\medskip
\noindent\textit{Editorial note.} The proof cites no bibliography item inside the counted span.\par\medskip
\noindent\textit{Allowed standard prerequisites (editorial).}\begin{itemize}
\item Numerical polynomial leading term and the Euler characteristic along a morphism. Stacks Project, Varieties, Lemmas lemma-numerical-polynomial-leading-term (tag 0BEN) and lemma-euler-characteristic-morphism (tag 0BEK). The first reduces the comparison to leading coefficients and the second gives the alternating-sum expansion of the Euler characteristic along f.
\item Dimension theory for locally algebraic schemes and generic finiteness. Stacks Project, Varieties, Lemmas lemma-dimension-locally-algebraic (tag 0A21) and lemma-dimension-fibres-locally-algebraic (tag 0B2L); Morphisms, Lemma morphisms-lemma-generically-finite (tag 02NW). Shows the generic point of Z is the unique point over the generic point of f(Z) and that f|\_Z is generically finite, which makes the degree available.
\item Vanishing of higher direct images of a finite morphism and the projection formula. Stacks Project, Cohomology of Schemes, Lemma coherent-lemma-finite-pushforward-coherent (tag 01Y6); Cohomology, Lemma cohomology-lemma-projection-formula (tag 01E8). Gives the vanishing of R\textasciicircum{}i f\_*O\_Z near xi and the rewriting of the twisted Euler characteristic.
\end{itemize}\medskip
\ifdefined\hypertarget\hypertarget{sourcepointer-1}{}\fi
\noindent\textit{Source reference.} Original reference 02NY: sources/stacks-chow/tags.txt, source lines 3497--3497. Cross-chapter standard prerequisite cited by the author inside the counted span; it is used as a pointer and its own proof is not part of this item.\par\medskip
\ifdefined\hypertarget\hypertarget{sourcepointer-2}{}\fi
\noindent\textit{Source reference.} Original lemma 0BEK: sources/stacks-chow/tags.txt, source lines 14353--14353. Cross-chapter standard prerequisite cited by the author inside the counted span; it is used as a pointer and its own proof is not part of this item.\par\medskip
\ifdefined\hypertarget\hypertarget{sourcepointer-3}{}\fi
\noindent\textit{Source reference.} Original reference 01E8: sources/stacks-chow/tags.txt, source lines 1825--1825. Cross-chapter standard prerequisite cited by the author inside the counted span; it is used as a pointer and its own proof is not part of this item.\par\medskip
\ifdefined\hypertarget\hypertarget{sourcepointer-4}{}\fi
\noindent\textit{Source reference.} Original lemma 0BEM: sources/stacks-chow/tags.txt, source lines 14355--14355. Cross-chapter standard prerequisite cited by the author inside the counted span; it is used as a pointer and its own proof is not part of this item.\par\medskip
\ifdefined\hypertarget\hypertarget{sourcepointer-5}{}\fi
\noindent\textit{Source reference.} Original lemma 0A21: sources/stacks-chow/tags.txt, source lines 12689--12689. Cross-chapter standard prerequisite cited by the author inside the counted span; it is used as a pointer and its own proof is not part of this item.\par\medskip
\ifdefined\hypertarget\hypertarget{sourcepointer-6}{}\fi
\noindent\textit{Source reference.} Original lemma 0B2L: sources/stacks-chow/tags.txt, source lines 13934--13934. Cross-chapter standard prerequisite cited by the author inside the counted span; it is used as a pointer and its own proof is not part of this item.\par\medskip
\ifdefined\hypertarget\hypertarget{sourcepointer-7}{}\fi
\noindent\textit{Source reference.} Original reference 02NW: sources/stacks-chow/tags.txt, source lines 3495--3495. Cross-chapter standard prerequisite cited by the author inside the counted span; it is used as a pointer and its own proof is not part of this item.\par\medskip
\ifdefined\hypertarget\hypertarget{sourcepointer-8}{}\fi
\noindent\textit{Source reference.} Original reference 01Y6: sources/stacks-chow/tags.txt, source lines 2545--2545. Cross-chapter standard prerequisite cited by the author inside the counted span; it is used as a pointer and its own proof is not part of this item.\par\medskip
\ifdefined\hypertarget\hypertarget{sourcepointer-9}{}\fi
\noindent\textit{Source reference.} Original lemma 0BEN: sources/stacks-chow/tags.txt, source lines 14356--14356. Cross-chapter standard prerequisite cited by the author inside the counted span; it is used as a pointer and its own proof is not part of this item.\par\medskip
\begin{lemma}
\label{lemma-intersection-number-and-pullback}
Let $k$ be a field. Let $f : Y \to X$ be a morphism of proper schemes over $k$.
Let $Z \subset Y$ be an integral closed subscheme of dimension $d$ and let
$\mathcal{L}_1, \ldots, \mathcal{L}_d$ be invertible $\mathcal{O}_X$-modules.
Then
$$
(f^*\mathcal{L}_1 \cdots f^*\mathcal{L}_d \cdot Z) =
\deg(f|_Z : Z \to f(Z)) (\mathcal{L}_1 \cdots \mathcal{L}_d \cdot f(Z))
$$
where $\deg(Z \to f(Z))$ is as in
Morphisms, Definition \ref{morphisms-definition-degree}
or $0$ if $\dim(f(Z)) < d$.
\end{lemma}

\begin{proof}
The left hand side is computed using the coefficient of $n_1 \ldots n_d$
in the function
$$
\chi(Y, \mathcal{O}_Z \otimes f^*\mathcal{L}_1^{\otimes n_1} \otimes
\ldots \otimes f^*\mathcal{L}_d^{\otimes n_d}) =
\sum (-1)^i
\chi(X, R^if_*\mathcal{O}_Z \otimes
\mathcal{L}_1^{\otimes n_1} \otimes \ldots \otimes
\mathcal{L}_d^{\otimes n_d})
$$
The equality follows from Lemma \ref{lemma-euler-characteristic-morphism}
and the projection formula
(Cohomology, Lemma \ref{cohomology-lemma-projection-formula}).
If $f(Z)$ has dimension $< d$, then the right hand side
is a polynomial of total degree $< d$ by
Lemma \ref{lemma-numerical-polynomial-from-euler}
and the result is true. Assume $\dim(f(Z)) = d$. Let
$\xi \in f(Z)$ be the generic point. By
dimension theory (see Lemmas \ref{lemma-dimension-locally-algebraic} and
\ref{lemma-dimension-fibres-locally-algebraic})
the generic point of $Z$ is the unique point of $Z$ mapping to $\xi$.
Then $f : Z \to f(Z)$ is finite over a nonempty open of $f(Z)$, see
Morphisms, Lemma \ref{morphisms-lemma-generically-finite}.
Thus $\deg(f : Z \to f(Z))$ is defined and in fact it is equal
to the length of the stalk of $f_*\mathcal{O}_Z$ at $\xi$
over $\mathcal{O}_{X, \xi}$. Moreover, the stalk of
$R^if_*\mathcal{O}_X$ at $\xi$ is zero for $i > 0$ because
we just saw that $f|_Z$ is finite in a neighbourhood of $\xi$
(so that Cohomology of Schemes, Lemma
\ref{coherent-lemma-finite-pushforward-coherent} gives the vanishing).
Thus the terms $\chi(X, R^if_*\mathcal{O}_Z \otimes
\mathcal{L}_1^{\otimes n_1} \otimes \ldots \otimes
\mathcal{L}_d^{\otimes n_d})$ with $i > 0$ have total
degree $< d$ and
$$
\chi(X, f_*\mathcal{O}_Z \otimes
\mathcal{L}_1^{\otimes n_1} \otimes \ldots \otimes
\mathcal{L}_d^{\otimes n_d})
=
\deg(f : Z \to f(Z)) \chi(f(Z),
\mathcal{L}_1^{\otimes n_1} \otimes \ldots \otimes
\mathcal{L}_d^{\otimes n_d}|_{f(Z)})
$$
modulo a polynomial of total degree $< d$ by
Lemma \ref{lemma-numerical-polynomial-leading-term}.
The desired result follows.
\end{proof}

\end{document}
